\section{作品概述}
\subsection{功能与特性}
系统将摄像头视频在本地FPGA上变换为绘画风格画面，主要数据通路不依赖主机推理。图像在DDR中缓存，轻量CNN完成空间纹理和颜色的联合变换，显示支持原图、风格图及同源左右对比。十种参数依次为梵高、浮世绘、剪纸、漫画、水墨、阿什维尔抽象、形体对比、蒙德里安、马赛克及雨夜厚涂。按键提供确定的现场控制；手势以个人模板、投票和四秒保持触发切换。UI显示采集、推理计数及当前风格，便于区分显示刷新与实际新图处理。

\subsection{应用领域}
作品面向互动艺术展示、数字画框、摄像头特效终端及嵌入式视觉教学。本地数据流有利于独立运行和控制视频传输开销，也提供从训练、定点导出、总线调度到FPGA显示的教学实例。手势仅作为交互入口，当前不承担身份识别等需要高可靠判决的任务。

\subsection{主要技术特点}
模型采用六通道DW/PW残差结构，在不下采样的条件下以膨胀卷积扩大邻域。硬件利用分条行缓存、DW三列bank及PW多上下文流水协调RAM和DSP开销。RISC-V与APB组成控制平面，128位AXI承担DDR像素流，避免CPU搬运整幅图像。完整帧发布、最终B响应屏障和安全消隐提交使采集、计算与显示可以并行。INT8权重、逐通道移位和冻结BN折叠将训练模型转换为可逐值核对的整数算术。

\subsection{主要性能指标}
\begin{table}[H]\centering\small
\begin{tabularx}{\linewidth}{@{}lX@{}}\toprule 指标 & 当前结果\\\midrule
图像空间尺寸 & 预处理及CNN为1280×720，显示720p60\\
风格数量 & 10组独立学习参数，共用单套加速器\\
资源与时序 & RAM230/256、DSP142/160；setup/hold +0.108/+0.026ns\\
处理速度 & CNN/DMA模型39.085fps；实际整机持续FPS待复测\\
数值一致性 & 历史50例完整720p、40例小尺寸软件核对；本轮新增20例QAT一致\\\bottomrule
\end{tabularx}
\end{table}

\Needspace{8\baselineskip}
\subsection{主要创新点}
\begin{enumerate}
\item 将全720p六通道网络与分条DW/PW流水共同设计，在有限RAM与DSP下保留逐像素推理。
\item 将帧所有权、参数身份、视图元数据和安全换帧统一管理，使多风格切换与同源对比一致。
\item 建立绘画纹理训练、整数参考、实际固件和CAD报告的对应链，分别评估画质、吞吐和显示刷新。
\end{enumerate}

\subsection{设计流程}
先在Xilinx平台以串口静态图验证协议、DDR和HDMI，再在Ti60平台接入SC431HAI与RISC-V。随后迭代到全分辨率六通道模型，经教师引导、绘画损失及整数QAT训练导出参数，完成RTL数值回归、布局布线与电气核对。每次变更保留旧版本，对新bit单独记录指纹。
